% GENERATED by scripts/make_proxy_tables.py -- do not edit by hand.
\begin{table}[t]
\centering
\caption{Ölçeklenebilirlik (stokastik navigasyon vekili, geodezik yürütme
modeli, 100 tohum ve sabit yoğunluk 5 görev/robot). Değerler senaryoların
ortalamasıdır. Uygunluk $\uparrow$, gecikme (ms) $\downarrow$. Her ölçekte en
iyi değer \textbf{kalın} gösterilmiştir. Kapalı çevrim Gazebo deneyleri 3, 5 ve
10 robotlu filoları kapsamaktadır.}
\label{tab:scalability}
\small
\setlength{\tabcolsep}{4pt}
\begin{tabular}{l cc cc cc }
\toprule
& \multicolumn{2}{c}{\textbf{3 robot}} & \multicolumn{2}{c}{\textbf{5 robot}} & \multicolumn{2}{c}{\textbf{10 robot}} \\
\cmidrule(lr){2-3}\cmidrule(lr){4-5}\cmidrule(lr){6-7}
\textbf{Yöntem} & Uyg. & Gec. & Uyg. & Gec. & Uyg. & Gec. \\
\midrule
\textbf{AHE-MRTA*} & \textbf{0.300} & 0.12 & \textbf{0.329} & 0.20 & \textbf{0.326} & 0.54 \\
BiG-MRTA & 0.267 & 0.02 & 0.293 & 0.05 & 0.298 & 0.17 \\
RoSTAM-EA & 0.249 & 1.93 & 0.261 & 3.48 & 0.245 & 7.00 \\
Cons-DBTA & 0.257 & \textbf{0.02} & 0.276 & \textbf{0.03} & 0.290 & \textbf{0.05} \\
\bottomrule
\end{tabular}
\end{table}
